% !TEX root = ../main.tex
\section{Propuesta del juego}\label{sec:propuesta}

\subsection{Descripción general}

El sistema es un juego de dominó por turnos, ejecutado en consola, en el que un jugador humano se enfrenta a entre uno y tres jugadores automáticos, en adelante CPU, en la misma máquina. Antes de comenzar, el jugador elige el tipo de dominó (doble~6 o doble~9), la cantidad de CPU, el nivel de dificultad de cada una (Principiante, Intermedio o Avanzado) y, de forma opcional, un límite de tiempo por turno. La partida se juega a una sola mano y siempre de forma individual: gana quien se quede primero sin fichas o, si el juego se cierra, quien conserve la menor suma de puntos.

El dominó concentra variación real en puntos identificables. El tipo de dominó modifica el tamaño del juego de fichas, el reparto y la cantidad de fichas disponibles para robar; los jugadores automáticos deciden su jugada con criterios distintos según su nivel; la mano atraviesa estados en los que se admiten jugadas diferentes; y las jugadas deben poder deshacerse, guardarse y recuperarse. Cada uno de estos puntos constituye una fuerza de diseño.

\subsection{Características del juego}

La Tabla~\ref{tab:criterios} relaciona las características del dominó que aportan material de diseño con su manifestación en el alcance definido.

\begin{table}[H]
\caption{Características del juego que aportan material de diseño}\label{tab:criterios}
\tablacompacta
\begin{tabularx}{\textwidth}{@{}L{4.4cm}Y@{}}
  \toprule
  \cab{Característica} & \cab{Manifestación en el dominó} \\ \midrule
  Estados de partida con reglas distintas & La partida pasa por los estados preparada, salida, en curso, cerrada y terminada. En la salida solo se admite la ficha que determina la regla de apertura y con el juego cerrado ya no se admiten jugadas. \\ \filasep
  Entidades con comportamiento diferente & El jugador humano y los jugadores CPU de tres niveles deciden su jugada de forma distinta. \\ \filasep
  Variantes & Tipo de dominó, número de CPU, dificultad de cada CPU y límite de tiempo por turno. \\ \filasep
  Jugadas registrables y reversibles & Historial de la mano, deshacer la última jugada del humano, guardar y cargar la partida. \\ \filasep
  Estructura contenedora & Mesa con dos extremos abiertos, fichas para robar y mano de cada jugador. \\
  \bottomrule
\end{tabularx}
\end{table}

\subsection{Aspectos susceptibles de variación}

En el juego varían el tipo de dominó, la forma en que cada jugador decide su jugada, el control de tiempo del turno y el medio en el que se guarda la partida. Adicionalmente, a futuro podrían variar la forma de presentación (consola o interfaz gráfica) y el tipo de jugador, por ejemplo un jugador remoto, sin que estos cambios deban afectar al dominio.
